% =============================================================================
% Chapter 4: Data Storage — ML Systems Lecture Slides (Volume II)
% =============================================================================
% IMAGES NEEDED (SVGs converted to PDF):
%   - storage-compute-gap.pdf     - workload-inversion.pdf
%   - storage-hierarchy-pyramid.pdf - pipelined-loading.pdf
%   - gds-path.pdf                - storage-economics.pdf
%   - checkpoint-staging.pdf
% =============================================================================
\documentclass[aspectratio=169, 12pt]{beamer}
\usepackage{../../assets/beamerthememlsys}

\mlsyssetup{
  volume       = {Volume II},
  chapter      = {Chapter 4},
  logo         = {../../assets/img/logo-mlsysbook.png},
  instlogo     = {../../assets/img/logo-harvard.png},
  chaptertitle = {Data Storage},
}

% --- Fonts ---
\usepackage{fontspec}
\setsansfont{Helvetica Neue}[
  BoldFont={Helvetica Neue Bold},
  ItalicFont={Helvetica Neue Italic},
  BoldItalicFont={Helvetica Neue Bold Italic},
]
\setmonofont{JetBrains Mono}[Scale=0.85]

% --- Packages ---
\usepackage{booktabs}
\usepackage{amsmath}

% --- Image paths ---
\graphicspath{
  {images/}
}

% --- Chapter-specific macros ---
\newcommand{\IOWall}{I/O Wall}
\newcommand{\Tsave}{$T_{\text{save}}$}
\newcommand{\BWreq}{$BW_{\text{required}}$}

% --- Helper: safe image include ---
\newcommand{\safeimg}[2][width=\textwidth,keepaspectratio]{%
  \IfFileExists{images/#2}{\includegraphics[#1]{#2}}{%
    \IfFileExists{#2}{\includegraphics[#1]{#2}}{%
      \fbox{\parbox[c][2.5cm][c]{0.85\linewidth}{\centering\footnotesize\textcolor{midgray}{[Missing image]}}}%
    }%
  }%
}

% --- Section count ---
\setcounter{mlsystotalsections}{8}

\title{Data Storage}
\author{Vijay Janapa Reddi}
\institute{Harvard University}
\date{}

\begin{document}

% =============================================================================
% TITLE SLIDE
% =============================================================================
\mlsystitle{Data Storage}{Feed the Accelerators or Waste Them}{cover_storage.png}

% =============================================================================
% LEARNING OBJECTIVES
% =============================================================================
\begin{frame}{Learning Objectives}
\note{
% -- LINK: Ch3 wired the fleet; this chapter feeds it with data
Students have compute, network, and fabric. Without storage, GPUs starve.

% -- NARRATE: ``This chapter is about the invisible bottleneck: storage.
Everyone optimizes compute and network, but storage stalls waste just as much.''

% -- ENGAGE: ``Has anyone profiled a training job and found the GPU waiting for data?''
Most have not --- that is why storage bottlenecks go undetected.

% -- FLEX: [CORE] Sets expectations.
IF SHORT: Read objectives quickly.
}

\footnotesize
\begin{enumerate}\setlength\itemsep{0pt}
  \item Analyze how ML workloads \textbf{invert storage assumptions} and identify the \textbf{\IOWall{}}
  \item Compare the six tiers of the \textbf{ML storage hierarchy} (HBM through Archive)
  \item Apply the \textbf{data pipeline throughput equation} to size storage bandwidth
  \item Explain how \textbf{pipelining and prefetching} hide storage latency
  \item Describe how \textbf{GPU Direct Storage} bypasses the CPU
  \item Evaluate storage \textbf{economics} including hidden costs (egress, idle GPUs)
  \item Design \textbf{checkpoint staging} strategies to minimize training pause time
\end{enumerate}

\end{frame}

\begin{frame}{Visual Language}
\note{
% -- NARRATE: Same color system. Point briefly and move on.
% -- FLEX: [OPTIONAL] Skip if students have seen this three times.
}

\small
Throughout this course, colors carry meaning:

\vspace{0.3cm}
\begin{columns}[T]
  \begin{column}{0.45\textwidth}
    \begin{mlsyscard}{computestroke}
    \textbf{Blue} --- Compute / Processing\\
    {\footnotesize GPU ops, forward/backward pass, inference}
    \end{mlsyscard}
    \vspace{0.15cm}
    \begin{mlsyscard}{datastroke}
    \textbf{Green} --- Data / Memory\\
    {\footnotesize Data flow, caches, healthy paths}
    \end{mlsyscard}
  \end{column}
  \begin{column}{0.45\textwidth}
    \begin{mlsyscard}{routingstroke}
    \textbf{Orange} --- Routing / Scheduling\\
    {\footnotesize Load balancers, batch windows}
    \end{mlsyscard}
    \vspace{0.15cm}
    \begin{mlsyscard}{errorstroke}
    \textbf{Red} --- Error / Cost / Bottleneck\\
    {\footnotesize Loss, decode phase, waste}
    \end{mlsyscard}
  \end{column}
\end{columns}
\end{frame}



% =============================================================================
\section{The Fuel Line}
% =============================================================================

\begin{frame}{The Engine Without Fuel}
\note{
% -- LINK: Ch2 built compute, Ch3 wired the network; now the 500x storage-compute gap
The engine and wires are ready. Without fuel, they are expensive sculpture.

% -- NARRATE: Point to the crimson card: ``H100 HBM: 3.35 TB/s. Single NVMe: 7 GB/s.
Gap: 500x.'' ANALOGY: ``A Formula 1 engine connected to a garden hose.''

% -- ENGAGE: ``If H100 consumes at 3.35 TB/s and NVMe delivers 7 GB/s, how many
NVMe drives to match?'' Expected: about 480. That is an entire rack of storage.

% -- WARN: Students assume NVMe is ``fast enough'' because 7 GB/s sounds fast.
Correct: 7 GB/s is 500x slower than HBM consumption rate.

% -- FLEX: [CORE] The 500x gap is the central number of this chapter.
}

\small
\begin{columns}[T]
  \begin{column}{0.55\textwidth}
    The fleet has:
    \begin{itemize}\setlength\itemsep{0pt}
      \item \textcolor{computestroke}{\textbf{Compute}}: PetaFLOPS per node (Ch.~2)
      \item \textcolor{routingstroke}{\textbf{Network}}: TB/s fabric (Ch.~3)
      \item \textcolor{datastroke}{\textbf{Storage}}: ???
    \end{itemize}

    \vspace{0.15cm}
    \begin{mlsyscard}{crimson}
    An engine without fuel is expensive sculpture.\\[0.1cm]
    {\footnotesize H100 HBM: \textbf{3.35 TB/s}\\
    Single NVMe: \textbf{7 GB/s}\\
    Gap: \textbf{$\sim$500$\times$}}
    \end{mlsyscard}
  \end{column}
  \begin{column}{0.42\textwidth}
    \safeimg[width=\textwidth,height=5.5cm,keepaspectratio]{storage-compute-gap.pdf}
  \end{column}
\end{columns}

\end{frame}

\begin{frame}{Running Example: 175B Model Training}
\note{
% -- LINK: 500x gap introduced the problem; running example gives concrete numbers
These numbers thread through every section of the chapter.

% -- NARRATE: Walk through the table. ``3 TB training data. 1,050 GB per checkpoint.
But look at total checkpoint volume: 4.5 PB over 30 days. That DWARFS the training
data by over 1,000x. The surprise: checkpoint I/O, not data loading, is the dominant
storage workload for frontier models.''

% -- WARN: Students underestimate checkpoint volume. They focus on training data size
and ignore that checkpoints are written every 10 minutes for 30 days.

% -- FLEX: [CORE] These numbers recur in every subsequent slide.
IF SHORT: State the 4.5 PB checkpoint number and the 1,000x ratio.
}

\small
\textbf{A 175B parameter language model, 30-day training run, 256 nodes:}

\vspace{0.1cm}
\renewcommand{\arraystretch}{1.15}
{\footnotesize
\begin{tabular}{@{}lr@{}}
  \toprule
  \textbf{Component} & \textbf{Volume} \\
  \midrule
  Training dataset (compressed)         & 3 TB \\
  Model weights (FP16)                  & 350 GB \\
  Optimizer state (FP32)                & 1,400 GB \\
  Single checkpoint (full)              & 1,050 GB \\
  All checkpoints (30 days, 10 min)     & $\sim$4.5 PB \\
  \bottomrule
\end{tabular}
}

\vspace{0.2cm}
\begin{mlsyscard}{errorstroke}
{\footnotesize \textbf{Surprise}: Checkpoint I/O ($\sim$4.5 PB) dwarfs training data reads (3 TB) by over 1,000$\times$. Checkpoint storage, not data loading, is the dominant storage workload for frontier models.}
\end{mlsyscard}

\end{frame}

% =============================================================================
\section{Workload Inversion}
% =============================================================================

\begin{frame}{ML Inverts Every Storage Assumption}
\note{
% -- LINK: Running example showed volumes; workload inversion shows WHY traditional storage fails
A database administrator moving to ML infrastructure would get everything wrong.

% -- NARRATE: ``Traditional storage optimizes for IOPS: many small random reads.
ML needs throughput: few massive sequential streams. Every assumption inverts.''
Point to the I/O Wall callout: ``When storage throughput cannot keep up, GPUs
idle regardless of their computational power.''

% -- ENGAGE: ``Why does RAID 5 waste bandwidth for ML?''
Expected: parity calculation wastes 15--25\% of bandwidth for data that is
immutable and already backed in object storage.

% -- WARN: Students with database backgrounds assume IOPS matters.
Correct: ML storage is throughput-limited (GB/s), not IOPS-limited.

% -- FLEX: [CORE] Workload inversion is the conceptual foundation for the chapter.
}

% --- Layout: FULL-WIDTH diagram ---
\centering
\safeimg[width=0.92\textwidth,height=4.8cm,keepaspectratio]{workload-inversion.pdf}

\vspace{0.1cm}
\mlsysalert{The I/O Wall:}{When storage throughput cannot deliver training data as fast as accelerators consume it, GPUs idle regardless of their computational power.}

\end{frame}

\begin{frame}{The Small File Problem}
\note{
% -- LINK: Workload inversion showed throughput matters; small files kill metadata first
Before bandwidth matters, metadata operations serialize and collapse.

% -- NARRATE: ``One file per sample means open(), stat(), close() per image.
Metadata server: 100K--300K ops/s. 10,000 workers saturate it in under 1 second.''
Point to the table: ``Aggregating into shards reduces metadata ops by 10,000--100,000x
and converts random reads to sequential streams.''

% -- ENGAGE: ``Why do WebDataset and TFRecord exist?''
Expected: to aggregate small files into large sequential shards.

% -- WARN: Students store datasets as individual files (one JPEG per image).
This works on a laptop but collapses the PFS at scale.

% -- FLEX: [CORE] The small file problem is the first bottleneck teams hit.
IF SHORT: State the metadata saturation number and the shard solution.
}

\small
\begin{columns}[T]
  \begin{column}{0.55\textwidth}
    \textbf{Storing one file per sample is catastrophic at scale:}

    \vspace{0.15cm}
    \begin{itemize}\setlength\itemsep{0pt}
      \item {\footnotesize Each file: \texttt{open()}, \texttt{stat()}, \texttt{close()}}
      \item {\footnotesize Metadata server: 100K--300K ops/s}
      \item {\footnotesize 10,000 workers $\to$ saturates in $<$1 sec}
    \end{itemize}

    \vspace{0.1cm}
    \begin{mlsyscard}{datastroke}
    \textbf{Solution}: Aggregate into large sequential shards (TFRecord, WebDataset, Parquet).\\[0.05cm]
    {\footnotesize Reduces metadata ops by 10,000--100,000$\times$}
    \end{mlsyscard}
  \end{column}
  \begin{column}{0.42\textwidth}
    \vspace{0.3cm}
    \renewcommand{\arraystretch}{1.2}
    {\scriptsize
    \begin{tabular}{@{}lrr@{}}
      \toprule
      & \textbf{Small Files} & \textbf{Shards} \\
      \midrule
      Files         & 14M      & 1,400 \\
      Metadata ops  & 42M/epoch & 4,200/epoch \\
      Read pattern  & Random   & Sequential \\
      Throughput    & 0.5 GB/s & 25 GB/s \\
      \bottomrule
    \end{tabular}
    }
  \end{column}
\end{columns}

\end{frame}

% =============================================================================
\section{Storage Hierarchy}
% =============================================================================

\begin{frame}{The Six-Tier ML Storage Hierarchy}
\note{
% -- LINK: Small file problem solved metadata; the hierarchy solves bandwidth cliffs
Students fixed metadata with shards. Now: how to bridge the 300,000x bandwidth gap.

% -- NARRATE: Walk bottom to top: ``HBM: 3.35 TB/s, 80 GB, expensive. DRAM: 200 GB/s.
NVMe: 7--25 GB/s. PFS: 4 GB/s/node. Object: 1 GB/s/node. Archive: tape.
Each tier is a prefetch buffer for the tier above.'' Emphasize: 300,000x span.

% -- ENGAGE: ``Which tier does the next batch need to be in BEFORE the current one finishes?''
Expected: HBM --- it is the only tier where computation happens.

% -- WARN: Students think of the hierarchy as static placement. Correct: data flows
dynamically through tiers, promoted on demand and demoted when idle.

% -- FLEX: [CORE] The six-tier hierarchy is the organizing framework for storage design.
}

% --- Layout: FULL-WIDTH pyramid ---
\centering
\safeimg[width=0.88\textwidth,height=4.8cm,keepaspectratio]{storage-hierarchy-pyramid.pdf}

\vspace{0.1cm}
\mlsysinsight{Design principle:}{Right data, right tier, right time. Each tier is a prefetch buffer for the tier above.}

\end{frame}

\begin{frame}{Bandwidth Cliffs Between Tiers}
\note{
% -- LINK: Pyramid showed the tiers; this table shows the cliffs between them
Students saw the hierarchy. Now: the specific bandwidth drops at each boundary.

% -- NARRATE: ``HBM to DRAM: 17x cliff. DRAM to NVMe: 15x. NVMe to PFS: 3--10x.
HBM empties in 24 ms but takes 400 ms to refill from DRAM. The next batch must
ALREADY be in HBM before the current one finishes. Every tier prefetches for the
tier above.''

% -- WARN: Students assume data can stream through the hierarchy in real time.
Correct: each cliff means latency hiding via prefetching, not streaming.

% -- FLEX: [CORE] The bandwidth cliffs motivate all pipelining and prefetching.
IF SHORT: State the 17x HBM-to-DRAM cliff and the prefetching requirement.
}

\footnotesize
\renewcommand{\arraystretch}{1.05}
{\scriptsize
\begin{tabular}{@{}llrl@{}}
  \toprule
  \textbf{Tier} & \textbf{Storage} & \textbf{Bandwidth} & \textbf{Cliff} \\
  \midrule
  0 & GPU HBM        & 3.35 TB/s       & \textcolor{errorstroke}{$\downarrow$ 17$\times$} \\
  1 & Host DRAM      & 200 GB/s        & \textcolor{errorstroke}{$\downarrow$ 15$\times$} \\
  2 & Local NVMe     & 7--25 GB/s      & \textcolor{errorstroke}{$\downarrow$ 3--10$\times$} \\
  3 & Parallel FS    & 4 GB/s/node     & \textcolor{errorstroke}{$\downarrow$ 10$\times$} \\
  4 & Object Storage & $\sim$1 GB/s/node & \\
  \bottomrule
\end{tabular}
}

\vspace{0.1cm}
\begin{mlsyscard}{errorstroke}
{\scriptsize \textbf{Implication}: HBM empties in 24 ms but takes 400 ms to refill from DRAM. The next batch must \emph{already} be in HBM before the current one finishes. Every tier prefetches for the tier above.}
\end{mlsyscard}

\end{frame}

\begin{frame}{Tier 0: GPU HBM --- The Destination}
\note{
% -- LINK: Bandwidth cliffs showed the flow; Tier 0 is the destination where computation happens
HBM is not just fast memory --- it is the ONLY place computation can occur.

% -- NARRATE: Walk through the table: ``175B model: weights, optimizer, activations,
gradient buffers, NCCL buffers. Total: 67--79 GB out of 80 GB. Less than 15\%
free for prefetch. Late batch = GPU stall. Early batch = no room.
Just-in-time, not just-too-late.''

% -- WARN: Students think 80 GB is plenty. Correct: 85--98\% is occupied by model state.
The prefetch buffer is razor-thin.

% -- FLEX: [CORE] HBM capacity constraints drive all downstream pipeline design.
IF SHORT: State the 85--98\% occupancy and the just-in-time requirement.
}

\small
\begin{columns}[T]
  \begin{column}{0.55\textwidth}
    \textbf{HBM is the only tier where weights and activations reside during computation.}

    \vspace{0.1cm}
    {\footnotesize
    \renewcommand{\arraystretch}{1.15}
    \begin{tabular}{@{}lr@{}}
      \toprule
      \textbf{Component (per GPU)} & \textbf{Size} \\
      \midrule
      Weights (FP16, tensor-parallel) & 43.75 GB \\
      Optimizer (ZeRO-3)              & 5.5 GB \\
      Activations                     & 10--20 GB \\
      Gradient buffers                & 5.5 GB \\
      NCCL buffers                    & 2--4 GB \\
      \midrule
      \textbf{Total occupied}         & \textbf{67--79 GB} \\
      \bottomrule
    \end{tabular}
    }
  \end{column}
  \begin{column}{0.42\textwidth}
    \vspace{0.1cm}
    \begin{mlsyscard}{crimson}
    \textbf{H100}: 80 GB HBM\\[0.1cm]
    {\footnotesize 85--98\% occupied.\\
    $<$15\% free for prefetch.\\[0.1cm]
    Late batch = GPU stall.\\
    Early batch = no room.\\[0.1cm]
    \textbf{Just-in-time, not just-too-late.}}
    \end{mlsyscard}
  \end{column}
\end{columns}

\end{frame}

% --- ACTIVE LEARNING 1: Predict ---
\begin{frame}{Predict: Text vs.\ Image Bandwidth}
\note{
% -- LINK: HBM showed tight capacity; now students estimate how much bandwidth different
modalities need from storage.

% -- NARRATE: ``60 seconds. Estimate the aggregate storage bandwidth for text training
vs image training. Same cluster, same iteration time.''
The 2,400x difference surprises most students. Ask 2--3 answers before revealing.

% -- ENGAGE: This IS the active learning moment. Do NOT reveal yet.
Listen for whether students account for the modality difference.

% -- FLEX: [CORE] This prediction creates the surprise that motivates tiered storage.
}

\centering
\vspace{0.8cm}
{\Large\bfseries Think--Estimate--Share}

\vspace{0.5cm}
{\large A 2,048-GPU cluster with 200 ms iteration time.\\[0.2cm]
\textbf{Estimate} the aggregate storage bandwidth needed for:\\
(a) \textcolor{datastroke}{\textbf{Text}} training (4,096 tokens/GPU, 4 bytes/token)\\
(b) \textcolor{computestroke}{\textbf{Image}} training (256 images/GPU, 150 KB/image)}

\vspace{0.5cm}
{\normalsize Write your estimates. \textcolor{midgray}{(60 seconds)}}

\end{frame}

\begin{frame}{Reveal: Text vs.\ Image Bandwidth}
\note{
% -- LINK: Students just estimated; now the reveal
% -- NARRATE: ``Text: 160 MB/s --- a single NAS node handles it. Image: 393 GB/s ---
you need a high-performance parallel file system. The 2,400x gap between modalities
drives hierarchy design. Text is volume-heavy but bandwidth-light. Image training
is the storage stress test.''

% -- WARN: Students assume all training workloads have similar storage needs.
Correct: LLM text training is trivial for storage; vision training is brutal.

% -- FLEX: [CORE] The modality gap is a key insight for storage planning.
}

\small
\renewcommand{\arraystretch}{1.2}
{\footnotesize
\begin{tabular}{@{}lrrl@{}}
  \toprule
  & \textbf{Text} & \textbf{Image} & \textbf{Ratio} \\
  \midrule
  Batch size/GPU    & 4,096 tokens $\times$ 4 B & 256 imgs $\times$ 150 KB & \\
  Per-GPU data      & 16 KB                      & 38.4 MB                  & 2,400$\times$ \\
  Aggregate BW      & \textbf{160 MB/s}          & \textbf{393 GB/s}        & 2,400$\times$ \\
  \midrule
  Storage needed    & Single NAS node            & High-perf PFS            & \\
  \bottomrule
\end{tabular}
}

\vspace{0.15cm}
\begin{columns}[T]
  \begin{column}{0.48\textwidth}
    \begin{mlsyscard}{datastroke}
    {\footnotesize \textbf{Text}: bandwidth-light, volume-heavy. Bottleneck = dataset size \& checkpointing.}
    \end{mlsyscard}
  \end{column}
  \begin{column}{0.48\textwidth}
    \begin{mlsyscard}{computestroke}
    {\footnotesize \textbf{Image}: bandwidth-heavy. Bottleneck = storage system's ability to feed accelerators.}
    \end{mlsyscard}
  \end{column}
\end{columns}

\end{frame}

% =============================================================================
\section{Pipeline Equation}
% =============================================================================

\mlsysfocus{The Data Pipeline Equation}{%
$BW_{\text{req}} = N_{GPUs} \times U_{target} \times \dfrac{S_{batch}}{T_{iteration}}$%
}

\begin{frame}{Pipeline Equation: Four Levers}
\note{
% -- LINK: Modality comparison showed bandwidth needs; the pipeline equation sizes them precisely
Students saw text vs image bandwidth. This equation generalizes the calculation.

% -- NARRATE: Walk through the worked example: ``256 GPUs, 80\% utilization, 256 images
at 150 KB per batch, 200 ms iteration. Result: 39.3 GB/s sustained, continuously,
for days. Miss this target and GPUs idle.''
Key insight: ``None of the four levers is free. Fewer GPUs costs throughput.
Lower utilization wastes hardware. Smaller batches hurt convergence.''

% -- WARN: Students tweak one variable without accounting for the others.
Correct: the equation shows all four are coupled.

% -- FLEX: [CORE] The pipeline equation is the diagnostic tool for storage sizing.
IF SHORT: Show the worked example result (39.3 GB/s), skip variable walkthrough.
}

\small
\begin{columns}[T]
  \begin{column}{0.48\textwidth}
    {\footnotesize
    \renewcommand{\arraystretch}{1.2}
    \begin{tabular}{@{}lp{4.5cm}@{}}
      \toprule
      \textbf{Variable} & \textbf{Meaning} \\
      \midrule
      $N_{GPUs}$ & Number of accelerators \\
      $U_{target}$ & Target utilization (0.8--0.95) \\
      $S_{batch}$ & Local batch size (bytes) \\
      $T_{iteration}$ & Forward-backward time \\
      \bottomrule
    \end{tabular}
    }

    \vspace{0.2cm}
    \mlsysconcept{Key:}{No lever is free. Reducing any input has a cost.}
  \end{column}
  \begin{column}{0.48\textwidth}
    \textbf{Worked example (ImageNet, 256 GPUs):}

    \vspace{0.1cm}
    {\footnotesize
    $256 \times 0.80 \times \frac{256 \times 150\text{KB}}{0.2\text{s}}$\\[0.2cm]
    $= 256 \times 0.80 \times 192\text{ MB/s}$\\[0.2cm]
    $= \textbf{39.3 GB/s}$
    }

    \vspace{0.15cm}
    \begin{mlsyscard}{errorstroke}
    {\footnotesize Must sustain \textbf{39.3 GB/s} continuously for days.\\
    Fall short $\to$ accelerators idle.}
    \end{mlsyscard}
  \end{column}
\end{columns}

\end{frame}

\begin{frame}{The Data Stall Ratio}
\note{
% -- LINK: Pipeline equation sized bandwidth; data stall ratio diagnoses if you hit it
The equation tells you what you NEED. The stall ratio tells you what you GET.

% -- NARRATE: ``Below 2\%: healthy. 2--10\%: marginal, will break at scale. Above 10\%:
clear storage bottleneck --- redesign the data path. Economic consequence:
5\% stall on 1,000 GPUs at \$2/GPU-hr for one week = \$16,800 wasted.''

% -- WARN: Students ignore stall ratios below 5\% as ``acceptable.'' Correct: at
scale, 5\% = thousands of dollars per week. Storage monitoring is economic, not optional.

% -- FLEX: [CORE] The stall ratio is the operational diagnostic for storage health.
IF SHORT: State the three thresholds and the \$16.8K/week number.
}

\small
$$\text{Data Stall \%} = \frac{T_{\text{step}} - T_{\text{compute}}}{T_{\text{step}}} \times 100$$

\vspace{0.15cm}
\renewcommand{\arraystretch}{1.15}
{\footnotesize
\begin{tabular}{@{}lll@{}}
  \toprule
  \textbf{Stall Ratio} & \textbf{Diagnosis} & \textbf{Action} \\
  \midrule
  $<$ 2\%   & Healthy pipeline            & Monitor, do not over-invest \\
  2--10\%   & Marginal (will break at scale) & Audit pipeline, add prefetch \\
  $>$ 10\%  & \textcolor{errorstroke}{Clear storage bottleneck}   & \textcolor{errorstroke}{Redesign data path} \\
  \bottomrule
\end{tabular}
}

\vspace{0.15cm}
\begin{mlsyscard}{errorstroke}
{\footnotesize \textbf{Economic consequence}: 5\% stall on 1,000 GPUs at \$2/GPU-hr for 1 week = \textbf{\$16,800 wasted}. Storage monitoring is a first-order economic concern.}
\end{mlsyscard}

\end{frame}

\begin{frame}{Pipelining and Prefetching}
\note{
% -- LINK: Stall ratio diagnosed the problem; pipelining is the solution
Students know they have a stall. Pipelining overlaps CPU loading with GPU compute.

% -- NARRATE: Walk through the diagram: ``While the GPU computes batch N, the CPU
loads batch N+1. The GPU never sees the load latency --- as long as loading finishes
before compute does.'' Point to the caveat: ``Pipelining hides AVERAGE latency.
Buffer depth must absorb P99 tail latency to prevent stalls at scale.''

% -- ENGAGE: ``What happens during a P99.9 spike?''
Expected: the buffer drains and the GPU stalls. Deeper buffers absorb more tail.

% -- WARN: Students size buffers for average latency. Correct: at 10,000 workers,
P99.9 events happen every few minutes. Size for tail, not average.

% -- FLEX: [CORE] Pipelining is the fundamental technique for hiding storage latency.
}

% --- Layout: FULL-WIDTH diagram ---
\centering
\safeimg[width=0.92\textwidth,height=4.5cm,keepaspectratio]{pipelined-loading.pdf}

\vspace{0.1cm}
\mlsysalert{Critical caveat:}{Pipelining hides \emph{average} latency. Buffer depth must absorb \textbf{P99 tail latency} to prevent stalls at scale.}

\end{frame}

% --- ACTIVE LEARNING: Micro-Retrieval Cue ---
\begin{frame}{Quick Check}
\note{
% -- NARRATE: ``What is the 500x gap, and what technique hides it?''
Answer: HBM 3.35 TB/s vs NVMe 7 GB/s. Pipelining and prefetching hide it.
% -- FLEX: [CORE] Micro-retrieval reinforces the central numbers.
}

\centering
\vspace{1.0cm}
{\Large\bfseries Quick Check}

\vspace{0.6cm}
{\large What is the 500$\times$ storage-compute gap,}\\
{\large and what technique hides it?}

\vspace{0.5cm}
{\normalsize\textcolor{midgray}{15 seconds --- then I will cold-call.}}

\end{frame}



% --- ACTIVE LEARNING 2: Exercise ---
\begin{frame}{Your Turn: Size the Prefetch Buffer}
\note{
% -- LINK: Pipelining introduced the technique; this exercise sizes the buffer
Students know pipelining hides latency. Now they calculate the required buffer depth.

% -- NARRATE: ``90 seconds. Calculate minimum prefetch depth for average and P99.''
After pause: ``Average: ceil(250/200) = 2 batches. P99: ceil(500/200) = 3 batches.
The P99 answer is the correct engineering answer. Engineer for tail, not average.''

% -- ENGAGE: This IS the active learning moment.
After solution: ``What if P99 is 1000 ms?'' Expected: 5 batches --- but that
may exceed the HBM prefetch budget from the Tier 0 slide.

% -- WARN: Students will compute for average and declare victory. The P99 answer
is the one that prevents production stalls.

% -- FLEX: [CORE] Buffer sizing exercise builds practical engineering intuition.
}

\small
\begin{columns}[T]
  \begin{column}{0.55\textwidth}
    {\normalsize\bfseries Size the Prefetch Buffer}

    \vspace{0.15cm}
    {\footnotesize A training node has:}
    \begin{itemize}\setlength\itemsep{0pt}
      \item {\footnotesize $T_{\text{compute}}$ = 200 ms per batch}
      \item {\footnotesize $T_{I/O,\text{avg}}$ = 250 ms (avg load time)}
      \item {\footnotesize $T_{I/O,\text{P99}}$ = 500 ms (tail latency)}
    \end{itemize}

    \vspace{0.05cm}
    {\footnotesize \textbf{Calculate} min prefetch depth (in batches) for (a) average and (b) P99 latency.}

    {\footnotesize\textcolor{midgray}{(90 seconds --- then compare)}}
  \end{column}
  \begin{column}{0.42\textwidth}
    \pause
    \begin{mlsyscard}{datastroke}
    \textbf{Solution:}\\[0.05cm]
    {\scriptsize
    Depth = $\lceil T_{I/O} / T_{\text{compute}} \rceil$\\[0.05cm]
    Average: $\lceil 250/200 \rceil$ = \textbf{2 batches}\\
    P99: $\lceil 500/200 \rceil$ = \textbf{3 batches}\\[0.05cm]
    \textcolor{datastroke}{\textbf{Engineer for P99, not average!}}
    }
    \end{mlsyscard}
  \end{column}
\end{columns}

\end{frame}

% =============================================================================
\section{GPU Direct Storage}
% =============================================================================

\begin{frame}{GPU Direct Storage: The CPU Bypass}
\note{
% -- LINK: Pipelining hid latency with CPU help; GDS eliminates the CPU from the path
Students saw CPU-based pipelining. GDS bypasses the CPU entirely.

% -- NARRATE: Point to the diagram: ``Traditional: 3 hops, CPU copies at every step.
GDS: 1 hop, direct DMA from NVMe to GPU memory via PCIe. The CPU is not involved.''
Point to the insight callout: ``GDS frees 6 of 8 CPU cores from data copying,
reallocating them to augmentation that improves model quality.''

% -- WARN: Students think GDS accelerates all storage. Correct: GDS only works
for local or RDMA-attached NVMe. It does not help object storage or network FS.

% -- FLEX: [CORE] GDS is the state-of-the-art for local data loading.
IF SHORT: State the 1-hop vs 3-hop comparison and the CPU core savings.
}

% --- Layout: FULL-WIDTH diagram ---
\centering
\safeimg[width=0.92\textwidth,height=4.8cm,keepaspectratio]{gds-path.pdf}

\vspace{0.1cm}
\mlsysinsight{CPU bypass dividend:}{GDS frees 6 of 8 CPU cores from data copying, reallocating them to augmentation that improves model quality.}

\end{frame}

\begin{frame}{GDS: Quantitative Impact}
\note{
% -- LINK: GDS diagram showed the path; this table quantifies the CPU savings
Students saw the data path. Now: specific numbers for a real workload.

% -- NARRATE: ``8 GPUs, 8,000 images/sec each, 150 KB per image. Traditional:
120 us per image, all 8 CPU cores consumed. GDS: 30 us, only 2 cores.
6 freed cores can run augmentation, improving both throughput and model quality.''

% -- WARN: Students overlook the limitation: GDS only accelerates local or
RDMA-attached NVMe. Object storage and network FS still go through the CPU.

% -- FLEX: [OPTIONAL] The quantitative detail reinforces GDS but is not examinable.
IF SHORT: State the 4x latency reduction and 6 freed cores.
}

\small
\textbf{Training node: 8 GPUs, 150 KB images, 8,000 images/s per GPU}

\vspace{0.1cm}
\renewcommand{\arraystretch}{1.2}
{\footnotesize
\begin{tabular}{@{}lrr@{}}
  \toprule
  & \textbf{Traditional} & \textbf{GDS} \\
  \midrule
  CPU time per image     & 120 $\mu$s         & 30 $\mu$s \\
  Total CPU load         & 7.68 core-seconds/s & 1.92 core-seconds/s \\
  CPU cores for I/O      & $\sim$8 cores       & $\sim$2 cores \\
  Cores freed for augmentation & 0              & \textbf{6 cores} \\
  \bottomrule
\end{tabular}
}

\vspace{0.15cm}
\begin{columns}[T]
  \begin{column}{0.48\textwidth}
    \begin{mlsyscard}{datastroke}
    {\footnotesize \textbf{4$\times$ latency reduction} per transfer (120 $\to$ 30 $\mu$s). No kernel buffer copies.}
    \end{mlsyscard}
  \end{column}
  \begin{column}{0.48\textwidth}
    \begin{mlsyscard}{errorstroke}
    {\footnotesize \textbf{Limitation}: GDS only accelerates \textbf{local or RDMA-attached NVMe}. Does not help object storage.}
    \end{mlsyscard}
  \end{column}
\end{columns}

\end{frame}

% =============================================================================
\section{Economics}
% =============================================================================

\begin{frame}{The Storage Cost Iceberg}
\note{
% -- LINK: GDS optimized the data path; economics shows the cost of getting it wrong
Students optimized throughput. Now: the dollar cost of storage decisions.

% -- NARRATE: ``The visible cost --- per-GB storage price --- is the tip of the iceberg.
Below the waterline: egress fees, IOPS charges, and most importantly, idle GPU cost.
For a 100 TB dataset, storage is only 7\% of total data delivery cost.''

% -- WARN: Students compare storage providers on per-GB price alone. Correct:
total cost includes egress, IOPS, and idle GPU time from stalls.

% -- FLEX: [CORE] The cost iceberg reframes storage as an economic decision.
IF SHORT: State the 7\% number and the hidden cost categories.
}

% --- Layout: FULL-WIDTH diagram ---
\centering
\safeimg[width=0.88\textwidth,height=4.5cm,keepaspectratio]{storage-economics.pdf}

\vspace{0.1cm}
\mlsysalert{Counterintuitive:}{The cheapest storage tier is not the cheapest choice. Total cost of data delivery includes egress + idle GPUs.}

\end{frame}

\begin{frame}{Worked Example: Stream vs.\ Cache}
\note{
% -- LINK: Cost iceberg introduced hidden costs; this worked example quantifies them
Students saw the iceberg. Now: a concrete stream-vs-cache calculation.

% -- NARRATE: Walk through both options: ``Stream from S3: \$373,800/yr (egress dominates).
NVMe cache: \$91,800/yr. Savings: \$282,000/yr --- enough for 3 additional H100 GPUs.
Break-even: 2 epochs. Multi-epoch training always benefits from local caching.''

% -- WARN: Students assume cloud storage is cheap because per-GB rates are low.
Correct: egress at \$0.09/GB compounds to 4x the storage cost over 20 epochs.

% -- FLEX: [CORE] The \$282K savings number makes storage optimization visceral.
IF SHORT: Show the bottom line comparison (\$374K vs \$92K).
}

\footnotesize
\textbf{50 TB vision dataset, 20 epochs, 4 training runs/year}

\vspace{0.1cm}
\begin{columns}[T]
  \begin{column}{0.48\textwidth}
    \textbf{Option A: Stream from S3}\\[0.1cm]
    {\scriptsize
    Storage: \$13,800/yr\\
    Egress: 50 TB $\times$ 20 $\times$ \$0.09/GB = \$90K/run\\
    4 runs/yr: \$360K egress\\[0.05cm]
    \textbf{Total: \$373,800/yr}
    }
  \end{column}
  \begin{column}{0.48\textwidth}
    \textbf{Option B: NVMe cache}\\[0.1cm]
    {\scriptsize
    S3 storage: \$13,800/yr\\
    NVMe capacity: \$60,000/yr\\
    Egress (stage 1$\times$/run): \$18,000/yr\\[0.05cm]
    \textbf{Total: \$91,800/yr}
    }
  \end{column}
\end{columns}

\vspace{0.15cm}
\begin{mlsyscard}{datastroke}
{\scriptsize \textbf{NVMe caching saves \$282,000/year} --- enough for 3 additional H100 GPUs. Break-even: 2 epochs of reading. Multi-epoch training always benefits from local caching.}
\end{mlsyscard}

\end{frame}

\begin{frame}{Tiering Strategy}
\note{
% -- LINK: Worked example showed cache savings; tiering generalizes the strategy
Students saw one cache decision. This formalizes the four-tier strategy.

% -- NARRATE: Walk through the table: ``Hot (NVMe): active training data. Warm (PFS):
shared datasets, recent checkpoints. Cold (S3): all organizational data. Archive (Glacier):
compliance and reproducibility.'' Emphasize: ``Automate with lifecycle policies, not
engineers. Data promotes when requested, demotes when idle.''

% -- WARN: Students manually place data and forget to demote. Result: hot tier overflows
while cold tier underutilized. Automate promotion and demotion.

% -- FLEX: [OPTIONAL] Tiering is a practical operations topic.
IF SHORT: Show the table, state the automation principle, move on.
}

\small
\renewcommand{\arraystretch}{1.15}
{\footnotesize
\begin{tabular}{@{}lllr@{}}
  \toprule
  \textbf{Tier} & \textbf{Storage} & \textbf{What Lives Here} & \textbf{\$/GB/mo} \\
  \midrule
  \textbf{Hot}     & Local NVMe + DRAM & Active training data, current batch    & \$0.10 \\
  \textbf{Warm}    & Parallel FS       & Shared datasets, recent checkpoints    & \$0.03 \\
  \textbf{Cold}    & Object (S3)       & All organizational data, versioned     & \$0.02 \\
  \textbf{Archive} & Glacier           & Compliance, reproducibility            & \$0.004 \\
  \bottomrule
\end{tabular}
}

\vspace{0.15cm}
\mlsysconcept{Lifecycle:}{Data promotes upward when a job requests it, demotes downward when idle. Automate with policies, not engineers.}

\vspace{0.1cm}
{\footnotesize \textbf{Staging patterns}: Pre-stage (best steady-state, slow start), On-demand (slow first epoch), Streaming (zero startup, no cache). Choice depends on epoch count and local NVMe capacity.}

\end{frame}

% --- ACTIVE LEARNING 3: Discussion ---
\begin{frame}{Discussion: Where Should Data Live?}
\note{
% -- LINK: Tiering introduced the four tiers; now students design a strategy
Students know the tiers. This forces them to reason about two different workloads.

% -- NARRATE: ``Turn to your neighbor. 70B language model on 2 TB text (single epoch)
AND a vision model on 50 TB images (20 epochs). Same cluster, same PFS.
Design the tiering strategy for each. 90 seconds.'' Cold-call 2--3 pairs.

% -- ENGAGE: This IS the active learning moment. No single right answer.
Key insight: single-epoch text may stream from object storage. Multi-epoch vision
MUST cache to NVMe. The same cluster needs different strategies simultaneously.

% -- WARN: Students apply one strategy to both workloads. Correct: access frequency
determines tier placement, and the two workloads have opposite patterns.

% -- FLEX: [CORE] This discussion synthesizes the entire storage chapter.
IF SHORT: Pose to whole class, take 2 answers, state the access-frequency principle.
}

\centering
\vspace{0.5cm}
{\large\bfseries Turn and Talk \textcolor{midgray}{(90 seconds)}}

\vspace{0.4cm}
{\normalsize Your team trains a 70B language model on 2 TB of text\\
(single epoch) and a vision model on 50 TB of images\\
(20 epochs). Same cluster, same PFS.\\[0.2cm]
\alert{Design the tiering strategy for each workload.}}

\vspace{0.4cm}
\begin{columns}[c]
  \begin{column}{0.22\textwidth}\centering\small Object\\Only\end{column}
  \begin{column}{0.22\textwidth}\centering\small PFS\\Cache\end{column}
  \begin{column}{0.22\textwidth}\centering\small NVMe\\Pre-stage\end{column}
  \begin{column}{0.22\textwidth}\centering\small Hybrid\\Strategy\end{column}
\end{columns}

\end{frame}

% =============================================================================
\section{Checkpoints}
% =============================================================================

\begin{frame}{Checkpoint Storage: The Burst Problem}
\note{
% -- LINK: Tiering covered read patterns; checkpoints are the dominant WRITE workload
Students optimized reading. Checkpoint writes are actually the bigger problem.

% -- NARRATE: ``All 256 nodes write simultaneously --- a checkpoint storm. 1,050 GB
hits the PFS in seconds. The PFS must be provisioned for PEAK burst, not average.''
Point to the staging concept: ``Write locally to NVMe (fast, non-blocking),
then replicate asynchronously to PFS in the background.''

% -- WARN: Students design checkpoint storage for average write bandwidth.
Correct: the burst is 10x steady-state. Under-provision and checkpointing
pauses training --- defeating its purpose.

% -- FLEX: [CORE] Checkpoint staging is a key architectural pattern for fault tolerance.
}

% --- Layout: FULL-WIDTH diagram ---
\centering
\safeimg[width=0.92\textwidth,height=4.5cm,keepaspectratio]{checkpoint-staging.pdf}

\vspace{0.1cm}
\mlsysconcept{Checkpoint storm:}{All nodes write simultaneously. PFS must be provisioned for \textbf{peak burst}, not average. Write locally, replicate asynchronously.}

\end{frame}

\begin{frame}{The Complete Data Path}
\note{
% -- LINK: Checkpoint staging showed one write pattern; this traces the complete data lifecycle
Students saw individual patterns. Now: the full lifecycle during a 30-day training run.

% -- NARRATE: Walk through each phase: ``Setup: stage 3 TB from object to NVMe once.
Training: NVMe to DRAM to HBM per batch. Checkpoint: HBM to NVMe (fast local write).
Replicate: NVMe to PFS (async background). Archive: PFS to object to Glacier.''
Emphasize: ``Checkpoint I/O (4.5 PB) dwarfs training data I/O (3 TB per epoch)
by 1,000x. For frontier models, checkpoint strategy has more impact on storage
design than training data pipeline optimization.''

% -- FLEX: [CORE] The complete data path closes the chapter's narrative arc.
IF SHORT: State the 4.5 PB vs 3 TB contrast and the checkpoint dominance insight.
}

\footnotesize
\textbf{Data lifecycle during a 30-day training run:}

\vspace{0.1cm}
\renewcommand{\arraystretch}{1.1}
{\scriptsize
\begin{tabular}{@{}lll@{}}
  \toprule
  \textbf{Phase} & \textbf{Data Path} & \textbf{Volume} \\
  \midrule
  \textbf{Setup}     & Object $\to$ NVMe (stage once)            & 3 TB \\
  \textbf{Training}  & NVMe $\to$ DRAM $\to$ HBM (per batch)     & 3 TB/epoch \\
  \textbf{Checkpoint} & HBM $\to$ NVMe (local, fast)              & 1,050 GB each \\
  \textbf{Replicate} & NVMe $\to$ PFS (async, background)         & 1,050 GB each \\
  \textbf{Archive}   & PFS $\to$ Object $\to$ Glacier (lifecycle) & $\sim$50 TB retained \\
  \bottomrule
\end{tabular}
}

\vspace{0.15cm}
\begin{mlsyscard}{crimson}
{\scriptsize \textbf{Counterintuitive}: Checkpoint I/O ($\sim$4.5 PB) dwarfs training data I/O (3 TB per epoch). For frontier models, \textbf{checkpoint staging strategy} has more impact on storage design than training data pipeline optimization.}
\end{mlsyscard}

\end{frame}

\begin{frame}{Checkpoint Coordination: Sync vs Async}
\note{
% -- LINK: Checkpoint staging showed the write path; this slide addresses WHEN
to checkpoint --- the coordination dimension.

% -- NARRATE: What to SAY while showing this slide
``Synchronous checkpointing: all workers pause, write state, resume. Simple but
expensive. 175B model checkpoint takes 60--120 seconds. At 10-minute intervals,
that is 10--20\% overhead. Asynchronous: write to local NVMe while training
continues. Background replication to PFS. Key risk: if a failure occurs during
async replication, the latest checkpoint may be incomplete. Solution: keep the
PREVIOUS checkpoint until replication confirms.''

% -- ENGAGE: ``What happens if async replication is slower than the checkpoint interval?''
Expected: checkpoints queue up and eventually overflow local NVMe. Must size
local storage for at least 2$\times$ checkpoint size.

% -- WARN: Students implement async without the two-checkpoint safety buffer.
If only one checkpoint exists and replication has not finished when failure
strikes, you lose ALL progress since the last completed checkpoint.

% -- FLEX: [CORE] Checkpoint coordination directly impacts training efficiency.
IF SHORT: State the 10--20\% overhead number and the async solution.
}

\small
\renewcommand{\arraystretch}{1.2}
{\footnotesize
\begin{tabular}{@{}llrl@{}}
  \toprule
  \textbf{Strategy} & \textbf{Impact on Training} & \textbf{Overhead} & \textbf{Recovery Risk} \\
  \midrule
  \textbf{Synchronous} & All workers pause & 10--20\% & None (consistent) \\
  \textbf{Async (local)} & Write NVMe, train continues & 1--3\% & Incomplete if NVMe fails \\
  \textbf{Async (staged)} & NVMe + background PFS & 2--5\% & \textcolor{datastroke}{\textbf{Safe}} (keep $n-1$) \\
  \bottomrule
\end{tabular}
}

\vspace{0.15cm}
\begin{columns}[T]
  \begin{column}{0.48\textwidth}
    \begin{mlsyscard}{errorstroke}
    {\scriptsize \textbf{Sync overhead}: 175B model, 60--120s per checkpoint, every 10 min = \textbf{10--20\% of training time} lost to I/O.}
    \end{mlsyscard}
  \end{column}
  \begin{column}{0.48\textwidth}
    \begin{mlsyscard}{datastroke}
    {\scriptsize \textbf{Async rule}: Always keep the \emph{previous} completed checkpoint on PFS until the current one is fully replicated.}
    \end{mlsyscard}
  \end{column}
\end{columns}

\end{frame}

% --- ACTIVE LEARNING: Prefetch Buffer Design Exercise ---
\begin{frame}{Your Turn: Design the Prefetch Buffer}
\note{
% -- LINK: Pipeline equation showed bandwidth requirements; this exercise applies it
to a concrete I/O hiding problem.

% -- NARRATE: ``90 seconds. Your training job reads 10 TB of images. NVMe delivers
7 GB/s, network delivers 25 Gbps (3.125 GB/s). Design the prefetch buffer to
hide I/O latency.'' After pause: ``GPU processes one batch in 0.5 s. Batch size:
200 MB. Required throughput: 200 MB / 0.5 s = 400 MB/s. NVMe delivers 7 GB/s ---
easily sufficient. Network at 3.125 GB/s --- also sufficient. But the question is
buffer size: you need enough to absorb P99 latency spikes. If P99 is 500 ms and
throughput is 400 MB/s, you need 200 MB buffer minimum. In practice: 2--4$\times$
margin = 400--800 MB buffer. For the full 10 TB dataset: stream from NVMe with
4 worker threads, prefetch 2$\times$ batch size ahead.''

% -- ENGAGE: This IS the active learning moment.
After: ``What if NVMe was only 1 GB/s (spinning disk)? Buffer = 400 MB/s $\times$ 3 s
latency = 1.2 GB. Slower storage demands larger buffers.''

% -- WARN: Students size the buffer for average latency, not P99. Tail latency
is what causes stalls --- the 500 ms spikes drain the buffer.

% -- FLEX: [CORE] I/O hiding is the most practical skill from this chapter.
}

\small
\begin{columns}[T]
  \begin{column}{0.55\textwidth}
    {\normalsize\bfseries Prefetch Buffer Design}

    \vspace{0.15cm}
    \textbf{Given:}
    \begin{itemize}\setlength\itemsep{0pt}
      \item {\footnotesize Training data: 10 TB of images}
      \item {\footnotesize NVMe throughput: 7 GB/s}
      \item {\footnotesize Network throughput: 25 Gbps (3.125 GB/s)}
      \item {\footnotesize Batch processing time: 0.5 s}
      \item {\footnotesize Batch size: 200 MB}
    \end{itemize}

    \vspace{0.1cm}
    \textbf{Design:} Buffer size to hide I/O latency.\\
    {\footnotesize\textcolor{midgray}{(90 seconds)}}
  \end{column}
  \begin{column}{0.42\textwidth}
    \pause
    \begin{mlsyscard}{datastroke}
    {\scriptsize \textbf{Solution:}\\
    Required: 200 MB / 0.5 s = 400 MB/s\\
    NVMe: 7 GB/s $\gg$ 400 MB/s \checkmark\\
    Network: 3.125 GB/s $\gg$ 400 MB/s \checkmark\\[0.1cm]
    \textbf{Buffer for P99 spikes:}\\
    P99 = 500 ms stall\\
    Buffer = 400 MB/s $\times$ 0.5 s\\
    = \textbf{200 MB minimum}\\
    Practice: 2--4$\times$ = \textbf{400--800 MB}\\[0.05cm]
    \textcolor{datastroke}{\textbf{Size for tail latency!}}}
    \end{mlsyscard}
  \end{column}
\end{columns}

\end{frame}

% =============================================================================
\section{Wrap-Up}
% =============================================================================

\begin{frame}{Fallacies}
\note{
% -- NARRATE: Four claims proven wrong with numbers. Spend extra time on the NVMe fallacy:
``7 GB/s sounds fast, but it is 500x slower than HBM. Without pipelining, every
batch transfer is a stall.'' The compression fallacy surprises: if CPU-bound,
decompression makes the bottleneck WORSE.
% -- FLEX: [CORE] IF SHORT: Cover the NVMe and compression fallacies.
}

\small
\textbf{Fallacy:} \textit{NVMe is fast enough to feed GPUs without pipelining.}\\
{\footnotesize NVMe: 7 GB/s. H100 HBM: 3.35 TB/s. Gap: $\sim$500$\times$. Without pipelining, every batch transfer is a stall.}

\vspace{0.1cm}
\textbf{Fallacy:} \textit{Object storage latency does not matter because we prefetch everything.}\\
{\footnotesize Prefetching hides \emph{average} latency, not tail. P99 spikes of 500+ ms drain the buffer and stall accelerators.}

\vspace{0.1cm}
\textbf{Fallacy:} \textit{Compression always helps.}\\
{\footnotesize Only if I/O-bound. In a CPU-bound pipeline, decompression makes the bottleneck \emph{worse}.}

\vspace{0.1cm}
\textbf{Fallacy:} \textit{RAID 5/6 is a safe default.}\\
{\footnotesize For ML, training data is immutable and backed in object storage. RAID 0 maximizes throughput; parity wastes 15--25\% bandwidth.}

\end{frame}

\begin{frame}{Pitfalls}
\note{
% -- NARRATE: Four operational mistakes. The egress pitfall is the most expensive:
``100 TB times 10 reads times \$0.09/GB = \$90,000 in egress, far exceeding
the \$27,600 annual storage cost. Stage locally.'' The accelerator upgrade pitfall
connects to Ch2: ``Upgrading A100 to H100 without upgrading storage INCREASES
the stall ratio.''
% -- FLEX: [CORE] IF SHORT: Cover the small files and egress pitfalls.
}

\footnotesize
\textbf{Pitfall:} \textit{Using millions of small files for training data.}\\
Metadata operations serialize on the PFS metadata server. At 10,000 workers, metadata saturates in $<$1 sec. Aggregate into shards.

\vspace{0.1cm}
\textbf{Pitfall:} \textit{Ignoring egress costs.}\\
100 TB $\times$ 10 reads $\times$ \$0.09/GB = \$90,000 in egress, far exceeding annual storage cost (\$27,600). Stage locally.

\vspace{0.1cm}
\textbf{Pitfall:} \textit{Designing checkpoint storage for average write performance.}\\
Checkpoints burst at 10$\times$ steady-state. PFS must be provisioned for peak, not average.

\vspace{0.1cm}
\textbf{Pitfall:} \textit{Assuming faster accelerators automatically improve throughput.}\\
Upgrading A100 $\to$ H100 without upgrading storage \emph{increases} the data stall ratio. Every accelerator upgrade requires a storage audit.

\end{frame}

% --- RETRIEVAL PRACTICE ---

% --- MUDDIEST POINT ---
\begin{frame}{Muddiest Point}
\note{
% -- NARRATE: ``Write the one concept you found most confusing. Anonymous. One sentence.''
% -- FLEX: [CORE] Address top 2--3 confusions in next lecture's opening.
}

\centering
\vspace{1.0cm}
{\Large\bfseries What was the \alert{muddiest point} today?}

\vspace{0.8cm}
{\normalsize Write down the concept you found \textbf{most confusing}.}

\vspace{0.5cm}
{\small\textcolor{midgray}{Anonymous. One sentence. Submit before you leave.}}

\end{frame}

\begin{frame}{What Were the Key Ideas?}
\note{
% -- NARRATE: ``Close your notes. Write four key concepts. 90 seconds.''
Walk around. Do NOT show Key Takeaways yet.
% -- FLEX: [CORE] Retrieval practice.
}

\centering
\vspace{1.5cm}
{\Large\bfseries Close your notes.}

\vspace{0.8cm}
{\large Write down the \textbf{4 most important concepts} from today.}

\vspace{0.8cm}
{\normalsize\textcolor{midgray}{90 seconds --- no peeking.}}

\end{frame}

\begin{frame}{Key Takeaways}
\note{
% -- LINK: Students wrote their own list; now compare.
% -- NARRATE: ``Did you capture: 500x gap, 300,000x hierarchy span, 4.5 PB checkpoints,
\$282K savings from NVMe caching?''
% -- FLEX: [CORE] Official summary. Read every bullet.
}

\scriptsize
\begin{itemize}\setlength\itemsep{0pt}
  \item \textbf{500$\times$ storage-compute gap}: HBM at 3.35 TB/s vs.\ NVMe at 7 GB/s. The \IOWall{} idles accelerators.
  \item \textbf{Six-tier hierarchy}: HBM $\to$ DRAM $\to$ NVMe $\to$ PFS $\to$ Object $\to$ Archive. 300,000$\times$ bandwidth span.
  \item \textbf{Workload inversion}: ML needs throughput (GB/s), not IOPS. Sequential streaming, RAID 0.
  \item \textbf{Pipeline equation}: \BWreq{} $= N_{GPUs} \times U_{target} \times S_{batch} / T_{iteration}$. Miss it $\to$ GPUs idle.
  \item \textbf{GPU Direct Storage}: Bypasses CPU, 4$\times$ latency reduction, frees cores for augmentation.
  \item \textbf{Economics}: Storage is 7\% of total data cost. Egress and idle GPUs dominate. NVMe cache saves \$282K/yr.
  \item \textbf{Checkpoint staging}: Write NVMe, replicate async. 175B model: $\sim$4.5 PB checkpoints in 30 days.
\end{itemize}

\end{frame}

\begin{frame}{References}
\note{
% -- NARRATE: ``Start with Mohan et al. for data stalls and NVIDIA GDS design guide.''
% -- FLEX: [OPTIONAL] Skip if running short.
}

\small
\mlsysref{Mohan+21}{Mohan et al. ``Analyzing and Mitigating Data Stalls in DNN Training.'' VLDB 2021.}
\mlsysref{NVIDIA-GDS}{NVIDIA. ``GPU Direct Storage Design Guide.'' Technical Report, 2023.}
\mlsysref{Rajbhandari+20}{Rajbhandari et al. ``ZeRO: Memory Optimizations Toward Training Trillion Parameter Models.'' SC 2020.}
\mlsysref{Eisenman+22}{Eisenman et al. ``Check-N-Run: a Checkpointing System for Training Deep Learning.'' NSDI 2022.}
\mlsysref{Audibert+23}{Audibert et al. ``tf.data: A Machine Learning Data Processing Framework.'' VLDB 2023.}

\end{frame}

\begin{frame}{Next Lecture: Distributed Training}
\note{
% -- LINK: Infrastructure trilogy is complete (compute, network, storage).
% -- NARRATE: ``The physical infrastructure is built. Compute, network, storage ---
all three layers of the Fleet Stack's foundation. Next: the algorithmic challenge.
How do you split a single training job across 1,000+ accelerators? Data parallelism,
tensor parallelism, pipeline parallelism, and expert parallelism.''
% -- FLEX: [CORE] Forward hook. IF SHORT: State the central question and dismiss.
}

\footnotesize
\begin{columns}[T]
  \begin{column}{0.48\textwidth}
    \vspace{0.3cm}
    \begin{mlsyscard}{crimson}
    {\normalsize\bfseries Infrastructure Complete}\\[0.2cm]
    {\footnotesize
    \textcolor{computestroke}{\textbf{Ch.~2}}: Compute (PetaFLOPS)\\
    \textcolor{routingstroke}{\textbf{Ch.~3}}: Network (TB/s fabric)\\
    \textcolor{datastroke}{\textbf{Ch.~4}}: Storage (tiered hierarchy)\\[0.15cm]
    \textbf{Next}: How to \emph{split} a single job across 1,000+ accelerators?}
    \end{mlsyscard}
  \end{column}
  \begin{column}{0.48\textwidth}
    \vspace{0.3cm}
    {\normalsize\textbf{Chapter 5: Distributed Training}}

    \vspace{0.15cm}
    \begin{itemize}\setlength\itemsep{2pt}
      \item \textcolor{datastroke}{\textbf{Data parallelism}} --- replicate model
      \item \textcolor{computestroke}{\textbf{Tensor parallelism}} --- split layers
      \item \textcolor{routingstroke}{\textbf{Pipeline parallelism}} --- split stages
      \item \textcolor{errorstroke}{\textbf{Expert parallelism}} --- route tokens
    \end{itemize}

    \vspace{0.1cm}
    {\footnotesize The algorithmic challenge that completes the fleet.}
  \end{column}
\end{columns}

\end{frame}



\appendix

\begin{frame}{Backup: Storage Equations Quick Reference}
\note{
% -- NARRATE: Backup --- key equations for problem sets.
% -- FLEX: [OPTIONAL] Only show if requested.
}

\footnotesize
\textbf{Key Storage Equations:}

\vspace{0.15cm}
\renewcommand{\arraystretch}{1.3}
\begin{tabular}{@{}ll@{}}
  \toprule
  \textbf{Metric} & \textbf{Formula} \\
  \midrule
  Pipeline BW required & $BW = N_{\text{GPUs}} \times U_{\text{target}} \times S_{\text{batch}} / T_{\text{iter}}$ \\
  Data stall ratio & $R_{\text{stall}} = \max(0, 1 - BW_{\text{storage}} / BW_{\text{required}})$ \\
  Checkpoint size & $C = N_{\text{GPUs}} \times (16M_{\text{params}} + A_{\text{activations}})$ \\
  Total checkpoint I/O & $V = C \times N_{\text{ckpts}} \times (1 + R_{\text{replicas}})$ \\
  Prefetch buffer & $B \geq \text{throughput} \times P99_{\text{latency}} \times 2$ \\
  \bottomrule
\end{tabular}

\vspace{0.15cm}
{\footnotesize \textbf{175B model checkpoint}: 1,050 GB. At 10-min intervals over 30 days: $\sim$4.5 PB total I/O.}

\end{frame}

\begin{frame}{Backup: Emerging Storage Technologies}
\note{
% -- NARRATE: Emerging technologies that may change the storage landscape.
% -- FLEX: [OPTIONAL] Forward-looking reference.
}

\footnotesize
\textbf{Technologies reshaping ML storage (2025--2030):}

\vspace{0.15cm}
\begin{itemize}\setlength\itemsep{2pt}
  \item \textbf{CXL Memory Pooling}: Disaggregated memory via CXL 3.0. Adds a tier between DRAM and NVMe at $\sim$50 GB/s. Could absorb checkpoint writes without NVMe.
  \item \textbf{Computational Storage}: NVMe drives with embedded compute for decompression, augmentation. Offloads CPU bottleneck in data pipelines.
  \item \textbf{GPU Direct Storage 2.0}: Direct NVMe-to-GPU DMA without CPU involvement. 4$\times$ latency reduction, frees CPU cores for augmentation.
  \item \textbf{Object Storage Tiering}: S3 Express One Zone: 10$\times$ lower latency than standard S3. May eliminate the need for PFS warm tier.
\end{itemize}

\end{frame}


\end{document}
